%% start of file `template.tex'.
%% Copyright 2006-2013 Xavier Danaux (xdanaux@gmail.com).
%
% This work may be distributed and/or modified under the
% conditions of the LaTeX Project Public License version 1.3c,
% available at http://www.latex-project.org/lppl/.


\documentclass[10pt,a4paper, roman]{moderncv}        % possible options include font size ('10pt', '11pt' and '12pt'), paper size ('a4paper', 'letterpaper', 'a5paper', 'legalpaper', 'executivepaper' and 'landscape') and font family ('sans' and 'roman')

% moderncv themes
\moderncvstyle{classic}                             % style options are 'casual' (default), 'classic', 'oldstyle' and 'banking'
\moderncvcolor{blue}                               % color options 'blue' (default), 'orange', 'green', 'red', 'purple', 'grey' and 'black'
%\renewcommand{\familydefault}{\sfdefault}         % to set the default font; use '\sfdefault' for the default sans serif font, '\rmdefault' for the default roman one, or any tex font name
%\nopagenumbers{}                                  % uncomment to suppress automatic page numbering for CVs longer than one page

% character encoding
\usepackage[utf8]{inputenc}                       % if you are not using xelatex ou lualatex, replace by the encoding you are using
\usepackage{fontawesome}

% adjust the page margins
\usepackage[scale=0.88]{geometry}
\setlength{\hintscolumnwidth}{2.1cm}                % if you want to change the width of the column with the dates
%\setlength{\makecvtitlenamewidth}{10cm}           % for the 'classic' style, if you want to force the width allocated to your name and avoid line breaks. be careful though, the length is normally calculated to avoid any overlap with your personal info; use this at your own typographical risks...

% personal data
\name{Lukas}{Burgholzer}
\title{Curriculum Vitae}                               % optional, remove / comment the line if not wanted
\phone[mobile]{+43~664~9754740}                   % optional, remove / comment the line if not wanted
%\phone[fixed]{+43~732~2468~9536}                    % optional, remove / comment the line if not wanted
\email{burgholzer@me.com}                               % optional, remove / comment the line if not wanted
\homepage{burgholzer.me}                         % optional, remove / comment the line if not wanted
\social[github]{burgholzer}
\social[linkedin][https://www.linkedin.com/in/lukas-burgholzer-7a1741a7/]{Lukas Burgholzer}
\extrainfo{{\small\upshape\faBirthdayCake}~22 January 1994}
\photo[64pt][0.4pt]{foto}                       % optional, remove / comment the line if not wanted; '64pt' is the height the picture must be resized to, 0.4pt is the thickness of the frame around it (put it to 0pt for no frame) and 'picture' is the name of the picture file

\usepackage{csquotes}
\usepackage[
	backend=biber,
	style=ieee,
	sortcites=true,
	sorting=ynt,
	url=true,
	eprint=true,
	giveninits=true,
	minnames=3,
	maxnames=6
]{biblatex}
\addbibresource{lit_header.bib}
\addbibresource{library.bib}
%\renewcommand*{\bibfont}{\normalfont\scriptsize}

\AtEveryBibitem{% Clean up the bibtex rather than editing it
	\clearname{editor}%
	\clearfield{series}%
	\clearfield{isbn}%
	\clearfield{issn}%
	\clearfield{day}%
	\clearfield{month}%
	\clearfield{place}%
	\clearfield{pages}%
	\clearfield{issue}%
	\clearfield{volume}%
%	\clearfield{doi}%
	\clearlist{location}%
	\clearfield{eventtitle}%
	\ifentrytype{inproceedings}{%
		\clearlist{publisher}
	}{}%
	\ifentrytype{software}{%
	}{%
		\clearfield{url}%
		\clearfield{urlyear}%
	}%
}

\pdfinfo{
	/Title (Curriculum Vitae - Lukas Burgholzer)
	/Author (Lukas Burgholzer)
	/Subject (Curriculum Vitae)
	/Keywords (CV, Résumé, Burgholzer)
}

\begin{document}

\makecvtitle

\section{Profile}
Quantum-software researcher, engineering leader, and CTO working across the quantum computing software stack and HPCQC integration.
I oversee the development of the \emph{Munich Quantum Toolkit (MQT)} and the \emph{Munich Quantum Software Stack (MQSS)}, including the Quantum Device Management Interface (QDMI), and co-founded MQSC.
As of November 2026, my work has accumulated over \textbf{3,000 citations} (\textbf{h-index 31}); MQT has \textbf{4.6 million PyPI downloads} and \textbf{1.4k GitHub stars}.

\section{Experience}
\cventry{2025--present}{CTO and Founding Member of MQSC}{Bootstrapped deep-tech startup that translates quantum-software research into industrial products}{\url{https://mq.sc}}{}{}
\cventry{2024--present}{Technical Lead of the Munich Quantum Software Stack (MQSS)}{Munich Quantum Valley initiative with 400+ researchers; lead developer and maintainer of QDMI, an emerging de facto standard for interoperable quantum-hardware integration in HPCQC software stacks}{\url{https://github.com/Munich-Quantum-Software-Stack/QDMI}}{}{}
\cventry{2023--present}{Chief Developer of the Munich Quantum Toolkit (MQT)}{Open-source ecosystem providing tools across the quantum computing software stack, including simulation, compilation, and verification; 4.6M PyPI downloads and 1.4k GitHub stars (November 2026)}{\url{https://mqt.readthedocs.io}}{}{}
\cventry{2023--present}{Senior Researcher}{Technical University of Munich}{Germany}{}{Chair for Design Automation. Mentorship of 8+ doctoral researchers; initiated and co-developed courses, adopted by DTU and Northwestern University.}
\cventry{2019--2023}{Doctoral Researcher}{Johannes Kepler University Linz}{Austria}{}{LIT Secure and Correct Systems Lab. Institute for Integrated Circuits.}
\cventry{2018--2019}{Graduate Research Associate}{MathConsult GmbH}{Linz, Austria}{}{}

\section{Education}
\cventry{2019--2023}{PhD Studies in Computer Science}{Johannes Kepler University Linz}{Austria}{}{Supervisor: Univ.-Prof. Dr. Robert Wille (Institute for Integrated Circuits), \textit{passed with distinction}}
\cventry{2017--2018}{Master's Degree in Industrial Mathematics}{Johannes Kepler University Linz}{Austria}{}{\textit{passed with distinction}}
\cventry{2016--2019}{Bachelor's Degree in Computer Science}{Johannes Kepler University Linz}{Austria}{}{\textit{passed with distinction}}  % arguments 3 to 6 can be left empty
\cventry{2013--2016}{Bachelor's Degree in Mathematics}{Johannes Kepler University Linz}{Austria}{}{\textit{passed with distinction}}

\section{Skills}
\cvitem{Languages}{German (native), English (fluent)}
\cvitem{Quantum software}{Tools across the quantum computing software stack: simulation, compilation, verification, optimization, mapping, and hybrid execution}
\cvitem{Algorithms}{Decision diagrams, automated and optimized compilation, and scalable verification}
\cvitem{Software engineering}{Modern C\texttt{++}, Python, Git, CMake, and CI/CD}
\cvitem{HPCQC and leadership}{HPCQC integration and QDMI; technical leadership for MQT, MQSS, and MQSC; mentorship of doctoral researchers}

\clearpage
\section{Awards and Selected Achievements}
\cvitem{2026}{EuroTech Rising Stars Award, Rising Innovation Stars category, EuroTech Universities Alliance.}
\cvitem{2026}{High-Tech Young Talent Award, Bavarian State Ministry of Science and the Arts (Bavarian Hightech Awards), for research with particularly high innovation potential and strong prospects for practical application.}
\cvitem{2025}{EDAA Outstanding Dissertation Award, European Design and Automation Association.}
\cvitem{2024}{ACM SIGDA Outstanding PhD Dissertation Award.}
\cvitem{2024}{GI Dissertation Award nomination, Gesellschaft für Informatik, Austrian Computer Society, and Swiss Informatics Society.}
\cvitem{2024}{Heinz Zemanek Prize, Austrian Computer Society.}
\cvitem{2024}{PhD Forum Best Poster Prize, Design, Automation and Test in Europe (DATE).}
\cvitem{2023}{Best Paper Award nomination, Asia and South Pacific Design Automation Conference.}
\cvitem{2022, 2023}{Mentor, NYUAD International Hackathon for Social Good in the Arab World (quantum computing).}
\cvitem{2022}{JKU Young Researcher's Award, Johannes Kepler University Linz.}
%\cvitem{2021}{Became IBM Certified Associate Developer for \enquote{Quantum Computation using Qiskit v0.2X}}
%\cvitem{2021}{Received a \enquote{Certificate of Quantum Excellence} from the 2021 Qiskit Global Summer School}
%\cvitem{2021}{Attained best possible score in the IBM Quantum Challenge 2021 ($>1400$ participants)}
%\cvitem{2021}{Top $40$ at the Xanadu QHack 2021 QML Challenge ($>400$ teams)}
%\cvitem{2020}{$5^\mathit{th}$ Place at the IBM Quantum Challenge Fall 2020 ($>1000$ participants)}
%\cvitem{2020}{Invitation to the IBM Qiskit Camp 2020, Yorktown Heights, NY, United States}
%\cvitem{2020}{Acceptance into the DAC Young Student Fellow Program}
%\cvitem{2019}{Invitation to the IBM Qiskit Camp Asia 2019, Tokyo, Japan}
%\cvitem{2019}{$3^\mathit{rd}$ Place at the IBM Quantum Challenge 2019}

\section{Research and Publications}
Quantum computers are rapidly advancing, yet the corresponding software lags behind. Through my research, I bridge this gap by transferring methods from classical design automation to the quantum realm---creating the essential software layer that makes quantum hardware usable.
My contributions include decision-diagram-based simulation methods that reduce workloads from days to minutes, automated and optimized compilation methods for diverse hardware platforms, compiler infrastructure based on MLIR, and efficient, scalable verification.
I oversee the development of the \emph{Munich Quantum Toolkit (MQT)} and the \emph{Munich Quantum Software Stack (MQSS)}, including QDMI---an emerging de facto standard for interoperable quantum-hardware integration---ensuring that powerful quantum processors can be integrated into high-performance computing centers and utilized by end-users.
My work covers the entire quantum computing software stack: from efficient data structures and simulation to compilation, verification, and hybrid execution.

I have published 80+ papers in leading venues including IEEE TCAD, ACM TQC, DAC, DATE, ICCAD, QCE, and SC. For a full list of publications, see \url{https://www.cda.cit.tum.de/team/burgholzer/}.

\section{Open-Source Software}
I am a strong advocate for open science. All methods developed in my research are released as part of the \emph{Munich Quantum Toolkit (MQT)} (\url{https://github.com/munich-quantum-toolkit}), forming a comprehensive, open-source ecosystem for quantum software.
These tools are engineered for performance (C++) and usability (Python), enabling thousands of researchers and developers worldwide to build quantum applications.
The MQT has become a staple in the community, with \textbf{4.6 million downloads} on PyPI and \textbf{1.4k stars} on GitHub (as of November 2026).
Tools natively integrate with IBM Qiskit, Xanadu's PennyLane, and \texttt{jeff}; they are deployed at HPC centers across Europe and beyond, including the Leibniz Supercomputing Centre (LRZ) and Oak Ridge National Laboratory, and serve as the technological backbone of MQSC.

For a full list of contributions, see \url{https://github.com/burgholzer}.

\nocite{
	burgholzerAdvancedEquivalenceChecking2021,burgholzerCharacteristicsReversibleCircuits2022,burgholzerEfficientConstructionFunctional2021,burgholzerHandlingNonunitariesQuantum2022,burgholzerHybridSchrodingerFeynmanSimulation2021,burgholzerLimitingSearchSpace2022,burgholzerImprovedDDbasedEquivalence2020,burgholzerPowerSimulationEquivalence2020,burgholzerQCECJKQTool2021,burgholzerRandomStimuliGeneration2021,burgholzerSimulationPathsQuantum2022,burgholzerVerifyingResultsIBM2020,burgholzerExploitingArbitraryPaths2022,burgholzerExploitingReversibleComputing2023,pehamEquivalenceCheckingParadigms2022,pehamEquivalenceCheckingParameterized2023,pehamEquivalenceCheckingQuantum2022,pehamOptimalSubarchitecturesQuantum2023,quetschlichAutomatedFrameworkRealizing2023,quetschlichCompilerOptimizationQuantum2023,schneiderSATEncodingOptimal2023,adarshSyReCSynthesizerMQT2022,berentSoftwareToolsDecoding2023,berentSATEncodingQuantum2022,willeBasisDesignTools2022,willeDecisionDiagramsQuantum2023,willeEfficientCorrectCompilation2020,willeJKQJKUTools2020,willeMQTQMAPEfficient2023,willeMappingQuantumCircuits2019,willeToolsQuantumComputing2022,willeVerificationQuantumCircuits2022,willeVisualizingDecisionDiagrams2021,hillmichReorderingDecisionDiagrams2022,grurlArraysVsDecision2020,burgholzerExploitingReversibleComputing2023,poggelRecommendingSolutionPaths2023, quetschlichMQTBenchBenchmarking2023,quetschlichReducingCompilationTime2023,quetschlichPredictingGoodQuantum2023,sanderHamiltonianSimulationDecision2023,pehamDepthoptimalSynthesisClifford2023,pehamOptimalSubarchitecturesQuantum2023,quetschlichHybridClassicalQuantum2023,burgholzerDesignAutomationTools2023,quetschlichApplicationawareQuantumCircuit2024,berentDecodingQuantumColor2024,burgholzerTensorNetworksDecision2023,quetschlichMQTPredictorAutomatic2024,sanderStrippingQuantumDecision2024,jiangFlatDDHighPerformanceQuantum2024,stadeAbstractModelEfficient2024,willeQDMIQuantumDevice2024,stadeOptimalStatePreparation2025,stadeTowardsSupportingQIR2025,matoMQTQuditsSoftware2024,prr_2025_1,mqt,burgholzer2025MQTCore,berentDecodingQuantumColor2024,burgholzer2025ForwardBackward,herzogJointCutting2025,kayaSoftwarePlatformSupport2024,qce_2025_1,qsw_2025_1,scws2025_1,date_2026_3,date_2026_2,sca_hpcasia_2026_1,sca_hpcasia_2026_2}

%\printbibliography
\end{document}


%% end of file `template.tex'.
